\documentclass[yuxi]{../template/Report}%方括号内写yuxi即生成预习报告
\settemplatedir{../template/}%设置模板路径

\exname{~~~~~~铁磁材料的磁滞回线和基本磁化曲线} %实验名称
\extable{} %实验桌号
\instructor{张建华} %指导教师
\class{信工2401} %班级
\name{姚舜瑜} %姓名
\stuid{3240100532} %学号

\nyear{2025} %年
\nmonth{11} %月
\nday{11} %日
\nweekday{二} %星期几，e.g. \nweekday{三}
\daypart{上午}%上午/下午

\redate{} %如有实验补做，补做日期
\resitu{} %情况说明：

\begin{document}
\maketitle%输出封面

\section{预习报告（10分）}

\subsection{实验综述（5分）}
\subsubsection{实验目的}
本实验主要研究铁磁材料的磁滞回线和基本磁化曲线，了解采用示波器测动态磁滞回线的原理，通过实验测量材料在不同磁场下的磁化强度，分析其磁性特征。

\subsubsection{实验原理}
铁磁性材料（例如铁、钴、镍 及其合金）是一种性能优异的磁性材料。在外磁场作用下能被强烈磁化，故在电工和电子技术中获得广泛应用。
当磁场强度为H的磁场作用于磁介质时，磁介质也要产生附加磁场，其磁化强度M与H的比值为磁化率$\chi_m=M/H$。此时磁介质中的总磁感应强度B为
\begin{equation}
    B = \mu_0(H+M) = \mu_0(1+\chi_m)H = \mu H
\end{equation}
式中$\mu_0$为真空磁导率，$\mu$为磁介质的磁导率。

对于铁磁质而言，M与H的关系非常复杂，它不仅取决于H，而且还与材料的磁化历史有关，即B-H曲线（或M-H曲线）不是单值函数，这种现象称为磁滞。以下这张图展示了典型的铁磁材料的磁滞回线和基本磁化曲线。

\begin{figure}[H]
    \centering
    \includegraphics[width=0.5\textwidth]{hysteresis_curve.png}
    \caption{铁磁材料的磁滞回线和基本磁化曲线示意图}
    \label{fig:hysteresis_curve}
\end{figure}

从没有磁化过的铁磁样品（H=0, B=0）开始，当磁场H从0开始增加时，磁感应强度B随之缓慢增加；当H增大到一定程度时，B随H的增加而急剧增大；当H继续增大时，B的增加又趋于缓慢；当H增大到Hs时，B达到饱和值Bs，此后即使H再增加，B也基本不变。这条Oa曲线称为基本磁化曲线。

如果在a点后减小磁场H，B将沿着另一条路径abcd下降。当H=0时，B不为0，而保留剩磁Br，这种现象称为磁滞。为消除剩磁，需加反向磁场，当反向磁场增大到-Hc时，B才为0，Hc称为矫顽力。继续增加反向磁场，样品将反向磁化，直到在-Hs时达到反向饱和。如果再减小反向磁场并加正向磁场，B将沿另一条路径defa变化，从而形成一个闭合的B-H曲线，称为磁滞回线。

磁滞现象的物理解释是铁磁体内存在许多被称为磁畴的小区域。每个磁畴内部，原子磁矩是平行排列的，因而具有很强的自发磁化。但在未经磁化的铁磁体内，各个磁畴的磁化方向各不相同，宏观上对外不显磁性。在外磁场作用下，那些磁化方向与外磁场方向相近的磁畴将扩大，而方向相反的磁畴将缩小，同时各磁畴的磁化方向也转向外磁场方向。当外磁场足够强时，所有磁畴都沿外磁场方向排列，达到磁饱和。当外磁场减弱时，由于磁畴壁的移动和磁矩的转向都有一定的阻力，磁畴结构不能完全恢复原状，从而产生剩磁。

\subsubsection{实验器材原理}
在本实验中，为了在示波器上观察磁滞回线，我们需要将磁场强度 H 和磁感应强度 B 转换为相互垂直的电压信号 $U_x$ 和 $U_y$。

\paragraph{H 的测量与转换}
我们通过在样品的原线圈中通以交变电流 $I_1$ 来产生交变磁场。原线圈中的磁场强度 H 与电流 $I_1$ 成正比：
\begin{equation}
    H = \frac{N_1 I_1}{L}
\end{equation}
其中 $N_1$ 是原线圈的匝数，L 是样品的有效磁路长度。
将一个小的取样电阻 $R_1$ 与原线圈串联，其两端的电压降 $U_x$ 为：
\begin{equation}
    U_x = I_1 R_1
\end{equation}
联立以上两式，可得 H 与 $U_x$ 的关系：
\begin{equation}
    \boxed{H = \frac{N_1}{L R_1} U_x}
\end{equation}

\paragraph{B 的测量与转换}
根据法拉第电磁感应定律，样品中磁感应强度的变化会在副线圈中产生感应电动势 $\mathcal{E}_2$：
\begin{equation}
    \mathcal{E}_2 = -N_2 \frac{d\Phi}{dt} = -N_2 S \frac{dB}{dt}
\end{equation}
其中 $N_2$ 是副线圈的匝数，S 是样品的截面积。
将副线圈连接到一个由电阻 $R_2$ 和电容 C 组成的串联 RC 积分电路。流过该电路的电流 $I_2$ 为：
\begin{equation}
    I_2 = \frac{\mathcal{E}_2}{Z} = \frac{\mathcal{E}_2}{R_2 + 1/(j\omega C)}
\end{equation}
当电路参数满足 $R_2 \gg 1/(\omega C)$ 时，可以近似认为 $I_2 \approx \mathcal{E}_2 / R_2$。
此时，电容器 C 两端的电压 $U_y$ 为：
\begin{equation}
    U_y = \frac{Q}{C} = \frac{1}{C} \int I_2 dt \approx \frac{1}{CR_2} \int \mathcal{E}_2 dt
\end{equation}
将 $\mathcal{E}_2$ 的表达式代入积分，可得：
\begin{equation}
    U_y \approx \frac{1}{CR_2} \int \left(-N_2 S \frac{dB}{dt}\right) dt = -\frac{N_2 S}{C R_2} B
\end{equation}
因此，磁感应强度 B 与积分电路输出的电压 $U_y$ 的关系式为：
\begin{equation}
    \boxed{B = -\frac{C R_2}{N_2 S} U_y}
\end{equation}

\paragraph{示波器显示}
综上所述，我们得到了 H 和 B 与电压 $U_x$ 和 $U_y$ 的线性关系。将 $U_x$ 信号输入示波器的 X 轴，将 $U_y$ 信号输入示波器的 Y 轴，就可以在屏幕上得到一条闭合的 B-H 关系曲线，即磁滞回线。

\subsubsection{实验步骤}
\begin{enumerate}
    \item 连接电路：将原线圈与交流电源和取样电阻串联，副线圈连接到 RC 积分电路，再将积分电路的输出连接到示波器的 Y 轴输入端，取样电阻两端的电压连接到示波器的 X 轴输入端。
    \item 调节交流励磁电压单调下降至零，进行退磁处理，并将光斑调节到屏幕中心。
    \item 按步长增加交流励磁电压幅值，测量基本磁化曲线，读出磁滞回线的端点坐标$(U_A, U_B)$。
    \item 保持励磁电压不变，利用光标功能测量磁滞回线上的各点坐标$(U_x, U_y)$。
    \item 分析数据，绘制基本磁化曲线，并计算材料的剩磁 Br 和矫顽力 Hc。
\end{enumerate}

\subsection{实验重点（3分）}
\begin{enumerate}
    \item 理解铁磁材料的磁滞现象及其物理机制。
    \item 掌握利用示波器测量磁滞回线的方法和技巧。
    \item 熟悉数据处理方法，能够绘制基本磁化曲线并计算相关磁性参数。
\end{enumerate}
\subsection{实验难点（2分）}
\begin{enumerate}
    \item 正确连接和调试电路，正确操作和读取示波器数据。
    \item 准确测量磁滞回线上的各点坐标，确保数据的可靠性。
    \item 运用Origin绘图工具，绘制本实验图像并进行数据分析。
\end{enumerate}
\begin{fullreportonly}
\section{原始数据（20分）}
（将有老师签名的“自备数据记录草稿纸”的扫描或手机拍摄图粘贴在下方，完整保留姓名，学号，教师签字和日期。）

\section{结果与分析（60分）}
\subsection{数据处理与结果（30分）}
（列出数据表格、选择适合的数据处理方法、写出测量或计算结果。）

\subsection{误差分析（20分）}
（运用测量误差、相对误差或不确定度等分析实验结果，写出完整的结果表达式，并分析误差原因。）

\subsection{实验探讨（10分）}
（对实验内容、现象和过程的小结，不超过100字。）

\section{思考题（10分）}
（解答教材或讲义或老师布置的思考题，请先写题干，再作答。）
\end{fullreportonly}
\insertnotes
\end{document}